\chapter{A ring and the water before fire}\label{connection:rings}
\regionline{The hole is only the beginning}
\atlasfiguretitle{A ring and the water before fire: paired comparison}
\comparisonplate{Obwarzanek}{Bagel}{Twisted strands and a smoother ring. Both meet hot water before the oven.}

A hole draws the eye so strongly that it can hide the rest of a bread. Obwarzanek, bagel, ftira and reikäleipä all invite us to look through their centres. Once the silhouette has made the introduction, the comparison needs to move outward. Is the dough twisted? Has it entered hot water? Will the ring be eaten soft, opened for lunch or kept for later?

In Kraków's obwarzanek, strands are twisted together before the dough is joined into a ring. The protected product specification describes the pre-bake hot-water treatment, as well as the permitted toppings and production setting. The twist is part of the identity. Looking at the surface shows a kind of construction that the smoother bagel beside it does not need to share.\citep{obwarspec}

\section{Two records, two kinds of story}
The obwarzanek specification points to a purchase for Queen Jadwiga dated 2 March 1394. The bagel chapter draws on a different record: YIVO's account of Jewish Kraków statutes from 1610, including bagels among gifts connected with childbirth. These are encounters with named foods in particular documents. The court purchase and the community rule give us different glimpses of bread's social life.\citep{obwarspec,yivobagel}

One record puts a bread in an account of expenditure. The other places it among people marking an occasion. Neither needs a dramatic invention scene to be interesting. Ordinary documentary detail can bring us close to the reasons a food was recognised and named. A ring may belong to commerce, celebration and routine at once.

The two recipes in that chapter also give the home baker a useful contrast. The obwarzanki receive a brief dip, while the bagels stay in the bath longer. Their formulas and shaping differ too. When the finished crusts are compared, the water treatment should be read alongside the entire recipe, rather than treated as the only variable in a controlled experiment.\citep{obwarhome,hamelbagel}

\section{A ring finds another city}
The Montreal and New York chapter follows bagels into named businesses and urban working lives. Fairmount's history begins with Isadore Shlafman's Montreal bakery in 1919, then follows the move to Fairmount Avenue in 1949. St-Viateur's account describes Myer Lewkowicz learning the trade after arriving in Canada, and Joe Morena learning Yiddish alongside bakery work. The shop becomes a place where a practice can continue through people who did not all begin with the same language or family background.\citep{fairmount08,stviateur08}

This is a documented route through makers, rather than an arrow drawn between vaguely similar breads. A family history and an apprenticeship explain how work reaches another generation. The bread has a recognisable form, but the city gives it addresses, customers and arguments about preference.

The practical comparison printed here uses a honey bath for the Montreal-style rings and a malted bath for the water bagels. Their size and dough differ, so the differences can be tasted without declaring a winner. A wide opening, a thick wall, a seed coating and a cut face each offer something to notice. The ring is small enough to inspect closely and familiar enough to provoke a conversation.

\section{The rings that never enter the bath}
Ftira keeps the hole and changes the work around it. Its sourdough ring goes into the oven without the bagel's water bath. Once baked, it can be opened for a Maltese lunch with fish, tomatoes, capers and other ingredients. The centre's absence does not prescribe the crust or the crumb.\citep{ftira,ftirafilling}

Finnish reikäleipä offers another answer. The flatter rye bread described in the Finnish account could be kept on poles beneath the ceiling. Its central opening participates in storage. Compare it with the Swedish crispbread story and the hole begins to look less like a decorative mark and more like part of how a household makes room for bread.\citep{finroot,duxknack}

These examples make the ring a particularly good travelling question. Form connects the first glance, then purpose separates the breads. A twisted street bread, a bakery bagel, a filled sourdough ring and a rye bread prepared for keeping all leave a space in the middle. What fills the story is the work around that space.

\section{Hands behind the breakfast counter}
New York's bagel story also includes organised work. Adeena Sussman's account of the trade follows Local 338, established in 1907 among Jewish immigrant bakers. Its meetings continued in Yiddish into the middle of the twentieth century. The union concerned itself with wages, hours and benefits alongside the skilled making of a familiar food.\citep{unionbagel08}

This gives the ring another kind of continuity. A bread can carry a craft, but the people doing that work also need conditions in which to practise it. Machinery and changes in the industry alter those conditions. The shape may remain recognisable on a breakfast plate while the work behind it has changed substantially.

Look back at the hand-rolling and apprenticeship in the Montreal accounts, then forward to your own small batch. Your fingers can learn something about joining a rope, opening the hole and moving a wet ring. That experience offers a point of attention when reading a baker's story. The shop's scale and working life remain larger than the home exercise, but the action is easier to picture once you have tried it.

The next time a ring is handed over a counter, notice its portion size, seed coating and the way it is offered. Those choices bring dough and customer together. They also remind us that the city's bread is sustained by repeated work, not simply by a reputation that travellers carry away.

\begin{breadnote}{Compare the water baths}
Choose the home bagels on page~58, or the Montreal-style rings and malted water bagels on pages~161 and~162. Follow each bath as written. Photograph a cut ring and record its wall thickness, crust and chew. Ftira on page~32 gives the comparison an oven-baked ring with another purpose.
\end{breadnote}
